\section{Introduction}
\label{sec:intro}

Fault diagnosis of rotating machinery, particularly rolling element bearings, is paramount for ensuring the safety, reliability, and economic efficiency of modern industrial systems. Traditional physical modeling of rolling bearings \cite{harris2001rolling,mcfadden1984model} provides foundational mechanical understanding. In recent years, deep learning has revolutionized fault diagnosis \cite{liu2018artificial}. Foundational architectures like Convolutional Neural Networks (CNNs) \cite{jia2016deep,wen2018new,zhang2018deep} and their transfer learning variants \cite{li2020deep} alongside traditional models \cite{zhao2019deep,lei2020applications} have achieved significant success under the assumption that comprehensive labeled data for all possible fault types and severities are available during the training phase. However, in real-world industrial environments, machines predominantly operate under normal conditions, and obtaining massive annotated data for rare or severe fault states is logistically prohibitive. This data scarcity fundamentally hampers the deployment of conventional data-driven diagnostic algorithms. 

To overcome this limitation, Zero-Shot Learning (ZSL) has emerged as a promising paradigm in intelligent fault diagnosis \cite{xian2018zero,pan2009survey}. ZSL aims to identify unseen fault categories—those with zero training instances—by transferring knowledge from seen classes using auxiliary knowledge, typically abstracted as semantic embeddings. In recent years, Generative Zero-Shot Fault Diagnosis (ZSFD) frameworks have gained considerable traction. By leveraging generative models such as Variational Autoencoders (VAEs) \cite{kingma2013auto}, Generative Adversarial Networks (GANs) \cite{goodfellow2014generative}, and more recently, Diffusion Models \cite{ho2020denoising}, these approaches synthesize virtual samples or features for unseen faults based on their targeted semantic descriptions, transforming the ZSL problem into a standard supervised classification task.

The rapid evolution of generative techniques has further expanded these boundaries. Notably, diffusion models have consistently outperformed VAEs and GANs in generating high-fidelity, diverse samples \cite{dhariwal2021diffusion,nichol2021improved}, scaling remarkably well to high-resolution data synthesis \cite{rombach2022high}. Concurrently, as industrial applications involve critical safety requirements, the importance of uncertainty quantification in deep learning has surged \cite{kendall2017uncertainties,abdar2021review}. Furthermore, there is a growing consensus that purely data-driven models lack physical realism, driving the emergence of physics-informed machine learning \cite{karniadakis2021physics}. 

Despite their empirical success, existing ZSFD methodologies suffer from two critical limitations. First, contemporary deep generative models function strictly through a data-driven lens, entirely oblivious to the physical mechanics of the rotating systems they model. This leads to the synthesized vibration signals frequently exhibiting bizarre frequency distributions or violating basic physical principles (e.g., resonance frequencies and defect kinematics), causing a severe domain shift between generated samples and actual unseen faults. Second, while current semantic embedding networks \cite{schonfeld2019cada} map unseen classes into a deterministic feature space, they entirely lack a mechanism to quantify measurement and diagnostic uncertainty. In high-stakes industrial applications, such as aerospace and power generation, providing a deterministic classification without an associated confidence interval critically undermines the model's trustworthiness and applicability.

This paper addresses these foundational gaps by proposing a Physics-Constrained Diffusion Model (PC-Diffusion) integrated with a Bayesian semantic embedding framework. We formulate a novel diffusion model architecture whose reverse denoising process is rigidly regularized by the inherent physical laws governing bearing dynamics. By embedding a Physics-Informed Neural Network (PINN) capable of capturing the differential equations of bearing vibrations, the generative process is anchored to physical reality. Furthermore, we construct a double-tier semantic description leveraging both expert-defined attributes and Large Language Model (LLM) embeddings. To quantify uncertainty, we employ a Bayesian approach mapping semantic vectors into probability distributions rather than point estimates via Variational Inference. 

The primary contributions of this paper are summarized as follows:
\begin{itemize}
    \item We propose PC-Diffusion, the first diffusion-based generative model for zero-shot fault diagnosis that inherently incorporates physical mechanics. By applying novel constraint loss functions (frequency-domain energy conservation and time-domain periodicity), the model ensures the synthetic signals rigidly adhere to machine dynamics. The effectiveness of these constraints is rigorously validated through frequency spectrum comparison and physical consistency analysis, demonstrating that the generated signals faithfully preserve the fault characteristic frequencies dictated by bearing kinematics.
    \item We introduce a Bayesian semantic embedding framework that fuses dual-level semantics. This allows for explicit uncertainty quantification during zero-shot inference, filling a crucial gap in industrial application reliability.
    \item We conduct extensive experiments on standard benchmark datasets (CWRU and XJTU-SY). The results demonstrate that PC-Diffusion not only significantly outperforms state-of-the-art ZSL diagnostic frameworks in recognizing unseen faults but also provides highly interpretable confidence bounded predictions.
\end{itemize}

The remainder of this article is organized as follows. Section \ref{sec:related} reviews related work. Section \ref{sec:method} thoroughly details the proposed methodology. Experimental setups are presented in Section \ref{sec:experiments}, followed by results in Section \ref{sec:results}. Lastly, conclusions are drawn in Section \ref{sec:conclusion}.